% Part: sets-functions-relations
% Chapter: relations
% Section: orders

\documentclass[../../../include/open-logic-section]{subfiles}

\begin{document}

\olfileid[fa-IR]{sfr}{rel}{ord}
\olsection{ترتیب‌ها}

\begin{explain}
بسیاری از مقایسه‌های ما متضمن آن‌اند که اشیایی را، از جهتی معین،
«کوچک‌تر از»، «برابر با» یا «بزرگ‌تر از» اشیای دیگر وصف کنیم. این
مقایسه‌ها با روابط \emph{ترتیب} سروکار دارند. اما روابط ترتیب گونه‌های
متفاوتی دارند. برای نمونه، برخی ایجاب می‌کنند که هر دو شیءِ دلخواه با
یکدیگر مقایسه‌پذیر باشند و برخی چنین ایجابی ندارند. برخی همانی را در بر می‌گیرند
(مانند~$\le$) و برخی آن را کنار می‌گذارند (مانند~$<$). داشتن یک
رده‌بندی در اینجا سودمند خواهد بود.
\end{explain}

\begin{defn}[پیش‌ترتیب]
رابطه‌ای که هم بازتابی و هم تراگذری باشد \emph{پیش‌ترتیب} نامیده
می‌شود.
\end{defn}

\begin{defn}[ترتیب جزئی]
پیش‌ترتیبی که پادمتقارن نیز باشد \emph{ترتیب جزئی} نامیده می‌شود.
\end{defn}

\begin{defn}[ترتیب خطی]\ollabel{def:linearorder}
ترتیب جزئی‌ای که همبند نیز باشد \emph{ترتیب تام} یا
\emph{ترتیب خطی} نامیده می‌شود.
\end{defn}

\begin{ex}
هر ترتیب خطی یک ترتیب جزئی نیز هست، و هر ترتیب جزئی یک پیش‌ترتیب نیز
هست، اما عکس این گزاره‌ها برقرار نیست. رابطهٔ جهان‌شمول روی~$A$ یک
پیش‌ترتیب است، زیرا بازتابی و تراگذری است. اما اگر $A$ بیش از یک
!!{element} داشته باشد، رابطهٔ جهان‌شمول پادمتقارن نیست و بنابراین
ترتیب جزئی نیست.
\end{ex}

\begin{ex}
رابطهٔ \emph{طولانی‌تر نبودن} $\preccurlyeq$ را روی~$\Bin^*$ در نظر
بگیرید: $x \preccurlyeq y$ اگر و تنها اگر $\len{x} \le \len{y}$. این
رابطه یک پیش‌ترتیب (بازتابی و تراگذری) و حتی همبند است، اما ترتیب جزئی
نیست، زیرا پادمتقارن نیست. برای نمونه، $01
\preccurlyeq 10$ و $10 \preccurlyeq 01$، اما $01 \neq 10$.
\end{ex}

\begin{ex}
یکی از ترتیب‌های جزئی مهم، رابطهٔ $\subseteq$ روی مجموعه‌ای از
مجموعه‌هاست. این رابطه در حالت کلی ترتیب خطی نیست، زیرا اگر $a \neq b$
و $\Pow{\{a, b\}} = \{\emptyset, \{a\}, \{b\}, \{a,b\}\}$ را در نظر
بگیریم، می‌بینیم که $\{a\} \nsubseteq \{b\}$ و
$\{a\} \neq \{b\}$ و $\{b\}
\nsubseteq \{a\}$.
\end{ex}

\begin{ex}
رابطهٔ \emph{بخش‌پذیری بی‌باقی‌مانده} ترتیب جزئی‌ای به دست می‌دهد که
ترتیب خطی نیست. برای اعداد صحیح $n$ و~$m$، $n \mid m$ را به این معنا
می‌نویسیم که $n$، $m$ را (بی‌باقی‌مانده) تقسیم می‌کند؛ یعنی این رابطه
برقرار است اگر و تنها اگر عدد صحیحی~$k$ وجود داشته باشد که $m = kn$. این رابطه روی~$\Nat$
ترتیبی جزئی است، اما ترتیبی خطی نیست: برای نمونه، $2 \nmid 3$ و نیز
$3 \nmid
2$. اگر بخش‌پذیری را رابطه‌ای روی $\Int$ در نظر بگیریم، فقط یک
پیش‌ترتیب است، زیرا پادمتقارن نیست: $1 \mid -1$ و $-1 \mid 1$، اما
$1 \neq -1$.
\end{ex}

\begin{ex}
رابطهٔ \emph{امتداد} روی مجموعه‌ای از دنباله‌های~$A^*$ چنین است:
$s \sqsubseteq s'$ اگر و تنها اگر $s = \emptyseq$ (دنبالهٔ تهی)،
$s = s'$، یا $s = \tuple{s_1, \dots, s_n}$ و $s' =
\tuple{s_1, \dots, s_n, s_{n+1}, \dots, s_m}$. اگر $s \sqsubseteq s'$
باشد، می‌گوییم $s$ یک \emph{قطعهٔ آغازینِ}~$s'$ است. رابطهٔ امتداد
روی $A^*$ ترتیبی جزئی است، اما ترتیبی خطی نیست؛ برای نمونه، اگر
$a \neq b$، آنگاه $ab \not\sqsubseteq ba$ و $ba \not\sqsubseteq ab$.
\end{ex}

\begin{defn}[ترتیب اکید]
\emph{ترتیب اکید} رابطه‌ای است که غیربازتابی، نامتقارن و تراگذری باشد.
\end{defn}

\begin{defn}[ترتیب خطی اکید]\ollabel{def:strictlinearorder}
ترتیب اکیدی که همبند نیز باشد \emph{ترتیب تام اکید} یا
\emph{ترتیب خطی اکید} نامیده می‌شود.
\end{defn}

\begin{ex}
$\le$ ترتیب خطیِ متناظر با ترتیب خطی اکید~$<$ است. $\subseteq$ ترتیب
جزئیِ متناظر با ترتیب اکید~$\subsetneq$ است.
\end{ex}

هر ترتیب اکید $R$ روی~$A$ را می‌توان با افزودن قطر $\Id{A}$، یعنی با
افزودن همهٔ زوج‌های~$\tuple{x,
x}$، به یک ترتیب جزئی تبدیل کرد. (این عمل \emph{بستار بازتابیِ}~$R$
نامیده می‌شود.) برعکس، با آغاز از یک ترتیب جزئی می‌توان از راه حذف
کردن~$\Id{A}$ یک ترتیب اکید به دست آورد. دو نتیجهٔ بعدی این مطلب را
دقیق می‌کنند.

\begin{prop}\ollabel{prop:stricttopartial}
اگر $R$ ترتیبی اکید روی~$A$ باشد، آنگاه $R^+ = R \cup \Id{A}$ ترتیبی
جزئی است. افزون بر این، اگر $R$ ترتیبی خطی اکید باشد، آنگاه $R^+$
ترتیبی خطی است.
\end{prop}

\begin{proof}
فرض کنید $R$ ترتیبی اکید باشد؛ یعنی $R \subseteq A^2$ و $R$
غیربازتابی، نامتقارن و تراگذری باشد. بگذارید $R^+ = R \cup \Id{A}$.
باید نشان دهیم که $R^+$ بازتابی، پادمتقارن و تراگذری است.

$R^+$ آشکارا بازتابی است، زیرا $\tuple{x, x} \in \Id{A} \subseteq
R^+$ برای هر $x \in A$.

برای نشان‌دادن پادمتقارن‌بودن $R^+$، به برهان خلف فرض کنید $R^+xy$ و
$R^+yx$، اما $x \neq y$. چون $\tuple{x,y} \in R \cup \Id{A}$، ولی
$\tuple{x, y} \notin \Id{A}$، باید $\tuple{x, y} \in R$، یعنی
$Rxy$. به همین ترتیب،~$Ryx$. اما این با فرض نامتقارن‌بودن $R$ تناقض
دارد.

برای اثبات تراگذری، فرض کنید $R^+xy$ و $R^+yz$. اگر هر دو
$\tuple{x, y} \in R$ و $\tuple{y,z} \in R$ برقرار باشند، آنگاه $\tuple{x, z} \in
R$، زیرا $R$~تراگذری است. در غیر این صورت، یا $\tuple{x, y} \in
\Id{A}$، یعنی $x = y$، یا $\tuple{y, z} \in \Id{A}$، یعنی $y = z$؛
دست‌کم یکی از این دو حالت برقرار است.
در حالت نخست، بنا بر فرض $R^+yz$ را داریم و $x = y$؛ پس $R^+xz$. حالت
دوم نیز مشابه است. در هر دو حالت، $R^+xz$؛ بنابراین $R^+$ نیز تراگذری
است.

دربارهٔ بند «افزون بر این»، فرض کنید $R$ همبند نیز باشد. پس برای هر
$x \neq y$، دست‌کم یکی از $Rxy$ و~$Ryx$ برقرار است؛ یعنی دست‌کم یکی از
$\tuple{x, y} \in R$ و $\tuple{y, x} \in R$ برقرار است. چون $R \subseteq R^+$، این امر دربارهٔ $R^+$
نیز همچنان برقرار است؛ پس $R^+$ نیز همبند است.
\end{proof}

\begin{prop}\ollabel{prop:partialtostrict}
اگر $R$ ترتیبی جزئی روی~$A$ باشد، آنگاه $R^- = R \setminus \Id{A}$
ترتیبی اکید است. افزون بر این، اگر $R$ ترتیبی خطی باشد، آنگاه $R^-$
ترتیبی خطی اکید است.
\end{prop}

\begin{proof}
اثبات این گزاره به‌عنوان تمرین واگذار می‌شود.
\end{proof}

\begin{prob}
برهانی برای \olref[sfr][rel][ord]{prop:partialtostrict} به دست دهید.
\end{prob}

نتیجهٔ سادهٔ زیر نشان می‌دهد که ترتیب‌های خطی اکید خاصیتی شبیه اصل
گسترش دارند:

\begin{prop}\ollabel{prop:extensionality-strictlinearorders}
اگر $<$ ترتیبی خطی اکید روی $A$ باشد، آنگاه:
\[
  (\forall a, b \in A)((\forall x \in A)(x < a \liff x < b) \lif a = b).
\]
\end{prop}

\begin{proof}
فرض کنید $(\forall x \in A)(x < a \liff x < b)$. اگر $a < b$، آنگاه
$a <
a$، که با غیربازتابی‌بودن $<$ تناقض دارد؛ پس $a \nless b$. درست به
همین ترتیب، $b \nless a$. پس $a = b$، زیرا $<$ همبند است.
\end{proof}

\end{document}
